\subsection{下中心列}

命题 5. 设 \(\mathfrak{g}\) 是一个李代数， \(\mathrm{P}\) 是 \(\mathfrak{g}\) 的一个子模。我们把子模 \(\mathrm{P}_n\) 定义在 \(\mathfrak{g}\) 中，公式为 \(\mathrm{P}_1 = \mathrm{P}\) 以及 \(\mathrm{P}_{n+1} = [\mathrm{P},\mathrm{P}_n]\) （对每个 \(n\geqslant 1\) ）。那么

\begin{enumerate}
\def\labelenumi{(\arabic{enumi})}
\setcounter{enumi}{17}
\tightlist
\item
\end{enumerate}

\[
[ \mathrm{P} _ {m}, \mathrm{P} _ {n} ] \subset \mathrm{P} _ {m + n},\tag{18}
\]

\[
\mathrm{P} _ {n} = \sum_ {p = 1} ^ {n - 1} \left[ \mathrm{P} _ {p}, \mathrm{P} _ {n - p} \right] for n \geqslant 2.\tag{19}
\]

我们对 \(m\) 作归纳来证明 (18)。情形 \(m = 1\) 是显然的。由 Jacobi 恒等式，

\[
[ [ \mathrm{P}, \mathrm{P} _ {m} ], \mathrm{P} _ {n} ] \subset [ \mathrm{P} _ {m}, [ \mathrm{P}, \mathrm{P} _ {n} ] ] + [ \mathrm{P}, [ \mathrm{P} _ {m}, \mathrm{P} _ {n} ] ],
\]

即

\[
[ \mathrm{P} _ {m + 1}, \mathrm{P} _ {n} ] \subset [ \mathrm{P} _ {m}, \mathrm{P} _ {n + 1} ] + [ \mathrm{P}, [ \mathrm{P} _ {m}, \mathrm{P} _ {n} ] ].
\]

归纳假设表明 \([\mathbf{P}_m, \mathbf{P}_{n+1}] \subset \mathbf{P}_{m+n+1}\) 且 \([\mathbf{P}_m, \mathbf{P}_n] \subset \mathbf{P}_{m+n}\) ，从而

\[
[ \mathrm{P} _ {m + 1}, \mathrm{P} _ {n} ] \subset \mathrm{P} _ {m + n + 1} + [ \mathrm{P}, \mathrm{P} _ {m + n} ] = \mathrm{P} _ {m + n + 1}.
\]

由公式 (18)， \(\mathrm{P}_n\supset \sum_{p = 1}^{n - 1}[\mathrm{P}_p,\mathrm{P}_{n - p}]\supset [\mathrm{P}_1,\mathrm{P}_{n - 1}] = \mathrm{P}_n\) ，由此得 (19)。

当取 \(\mathbf{P} = \mathfrak{g}\) 时，序列 \((\mathrm{P}_n)\) 就是下中心列 \((\mathcal{C}^n\mathfrak{g})\) （对 \(\mathfrak{g}\) 而言）(第 I 章， § 1， no. 5)。于是：

命题 6. 设 g 是一个李代数， \((\mathcal{C}^{n}\mathfrak{g})_{n\geqslant 1}\) 是 g 的下中心列。那么

\[
[ \mathcal {C} ^ {m} \mathfrak {g}, \mathcal {C} ^ {n} \mathfrak {g} ] \subset \mathcal {C} ^ {m + n} \mathfrak {g} \quad for m \geqslant 1 and n \geqslant 1.
\]

推广第 I 章 § 4 no. 1 中的定义 1，我们称李代数 g 是幂零的，如果 \(\mathcal{C}^{n}\mathfrak{g} = \{0\}\) （对充分大的 \(n\) ）。幂零李代数 g 的幂零类是最小的整数 \(n\) ，它使得 \(\mathcal{C}^{n + 1}\mathfrak{g} = \{0\}\) 成立。

命题 7. 设 X 是一个集合，n 是一个 ≥1 的整数。

\begin{enumerate}
\def\labelenumi{(\alph{enumi})}
\item
  \(\mathrm{L}^{n+1}(\mathrm{X}) = [\mathrm{L}^{1}(\mathrm{X}), \mathrm{L}^{n}(\mathrm{X})]\) 。
\item
  模 \(L^{n}(X)\) 由形如 \([x_{1}, [x_{2}, \ldots, [x_{n-1}, x_{n}] \ldots]]\) 的元素生成，其中 \((x_{1}, \ldots, x_{n})\) 取遍由 X 的 n 个元素组成的序列的集合。
\item
  L(X) 的下中心列由 \(\mathcal{C}^n (\mathrm{L}(\mathrm{X})) = \sum_{p\geqslant n}\mathrm{L}^p (\mathrm{X})\) 给出
\item
  我们取 \(\mathfrak{g} = \mathrm{L}(\mathrm{X})\) 与 \(\mathrm{P} = \mathrm{L}^{1}(\mathrm{X})\) 应用命题 5。对 \(n\) 作归纳，由 (12) (no. 6) 与 (19) 推出等式 \(\mathrm{P}_{n} = \mathrm{L}^{n}(\mathrm{X})\) 。所欲证的关系式于是等价于定义 \([\mathrm{P}, \mathrm{P}_{n}] = \mathrm{P}_{n + 1}\) 。
\item
  这由 (a) 对 \(n\) 作归纳得到。
\item
  设 \(\mathfrak{g} = \mathrm{L}(\mathrm{X})\) 与 \(\mathfrak{g}_{n} = \sum_{p\geqslant n}\mathrm{L}^{p}(\mathrm{X})\) 。那么 \(\mathfrak{g} = \mathfrak{g}_1\) ，并且 no. 6 的公式 (13) 蕴涵 \([\mathfrak{g}_{n},\mathfrak{g}_{m}]\subset \mathfrak{g}_{n + m}\) ，特别地有 \([\mathfrak{g},\mathfrak{g}_n]\subset \mathfrak{g}_{n + 1}\) 。对 \(n,\mathcal{C}^{n}\mathfrak{g}\subset \mathfrak{g}_{n}\) 作归纳。另一方面，由 (a) 我们归纳地导出 \(\mathrm{L}^{n}(\mathrm{X})\subset \mathcal{C}^{n}\mathfrak{g}\) （对 \(n\) 归纳）。由于 \(\mathcal{C}^{n}\mathfrak{g}\) 是 \(\mathfrak{g}\) 的一个理想，关系式 \(\mathrm{L}^p (\mathrm{X})\subset \mathcal{C}^{n}\mathfrak{g}\) 蕴涵
\end{enumerate}

\[
\mathrm{L} ^ {p + 1} (\mathrm{X}) = [ \mathrm{L} ^ {1} (\mathrm{X}), \mathrm{L} ^ {p} (\mathrm{X}) ] \subset \mathcal {C} ^ {n} \mathfrak {g}
\]

这由 (a) 得出。因此有 \(\mathbf{L}^p (\mathbf{X})\subset \mathcal{C}^{n}\mathfrak{g}\) （对 \(p\geqslant n\) ），从而 \(\mathfrak{g}_n\subset \mathcal{C}^{n}\mathfrak{g}\)

推论. 设 \(\mathfrak{g}\) 是一个李代数， \((x_{i})_{i\in I}\) 是 \(\mathfrak{g}\) 的一个生成族。下中心列的第 \(n\) 项 \(\mathcal{C}^{n}\mathfrak{g}\) （对 \(\mathfrak{g}\) 而言）是由迭代括号 \([x_{i_1}, [x_{i_2}, \ldots, [x_{i_{p-1}}, x_{i_p}] \ldots]]\) 生成的模，其中 \(p \geqslant n\) ，且 \(i_1, \ldots, i_p\) 属于 I。

设 \(f\) 是 L(I) 到 g 中的同态，使得 \(f(i) = x_i\) （对一切 \(i \in I\) ）。由于 \((x_i)_{i \in I}\) 生成 g，有 g = f(L(I))，从而由第 I 章 § 1 no. 5 的命题 4 得 \(\mathcal{C}^{\bullet}g = f(\mathcal{C}^{\bullet}(L(I)))\) 。推论于是由命题 7 的断言 (b) 与 (c) 得出。

